\documentclass[titlepage,a4paper]{jsarticle}
\usepackage{../../../sty/import}% 各種パッケージインポート
\usepackage{../../../sty/title}% タイトルページの変更
\usepackage{../../../sty/answergrid}
%% タイトルページの変数
% レポートタイトル
\title{心理学特論第15回目出席票}
% 提出日
\expdate{\today}
% 科目名
\subject{心理学特論}
% 分野
\class{情報経営システム工学分野}
% 学年
\grade{M1}
% 学籍番号
\mynumber{24336488}
% 記述者
\author{本間三暉}

\begin{document}
% titleページ作成
\maketitle
\section*{keyword\#1}
ジャネーの法則
\section*{keyword\#2}
パッシブレスト
\section*{☆心理学特論の第15回テーマ「時間の心理学」について，簡潔にまとめてみましょう！
今回が今期「心理学特論」最後の授業です．授業アンケートへの回答も各自でお願いします．}
\section{「気がついたら時間がなくなっていた！」と思えた経験があれば，それをなるべくわかりやすく説明して下さい．
（何でそうなってしまったのか？という原因分析も含めて）}
私は，プログラミングをしているときに，「気がついたら何時間も経っていた」と感じることがよくある．
新しい機能を実装したり，不具合の原因を調べたりしていると，一つ解決するたびに次の改善点が見つかるため区切りを付けるタイミングを逃してしまう．
その結果，「少しだけ作業しよう」と思って始めても，気付けば3～4時間ほど経っていることがある．

このようになってしまう原因は，開発そのものが楽しく目標が明確であるため，高い集中状態になっているからだと考える．
また，「あと少しで完成する」「この不具合だけ直したい」という気持ちが続くことで，作業をやめるきっかけを失ってしまうことも大きな要因である．
そのため，時間の経過を意識しにくくなり，気付いたときには予定していた時間を大幅に超えていることが多い．

\section{「時短・効率化」のために，自分なりに考えて行動して上手くいった経験があれば，具体的にどの程度まで時間の短縮が図れたか，作業の効率化が達成できたのか，その工夫や心がけたことをなるべくわかりやすく説明して下さい．
（経験がない場合は，どうすればできるか？という目論見でも結構です）}
私は，レポートやスライドの作成にMicrosoft Officeを使うことをやめ，Markdownなどのコードベースで管理するようにした．
以前は，WordやPowerPointでレイアウトを調整するためにGUIを何度も操作する必要があり，少し文章を修正しただけで全体の見た目が崩れることがあった．
そのため，内容よりも体裁の修正に多くの時間を費やしていた．

コードベースで管理するようになってからは，内容の修正に集中できるようになり，レイアウトも自動で反映されるため，資料作成にかかる時間を短縮できた．
また，バージョン管理ツールを使っているため，ファイルを探したり最新版を確認したりする手間もなくなった．
以前は，提出したレポートをダウンロードフォルダから保存先へ移動し忘れ，定期的な整理の際に誤って削除してしまうこともあったが，現在はそのような管理ミスもなくなった．

このように，作業方法そのものを見直すことで，資料作成やファイル管理に費やす時間を減らし，本来取り組むべき内容の検討や改善に多くの時間を使えるようになった．

\end{document}